% ICASSP 2026 Conference Paper (polished with SRMR-consistent reward)
% Two-column format with 5-6 page target (includes references)

\documentclass[10pt,conference]{IEEEtran}
\IEEEoverridecommandlockouts

% Packages
\usepackage{cite}
\usepackage{amsmath,amssymb,amsfonts}
\usepackage{algorithmic}
\usepackage{algorithm}
\usepackage{graphicx}
\usepackage{textcomp}
\usepackage{xcolor}
\usepackage{booktabs}
\usepackage{multirow}
\usepackage{url}
\usepackage{float}
\usepackage{placeins}

\usepackage{hyperref}
\hypersetup{
    colorlinks=true,
    linkcolor=blue,
    citecolor=blue,
    urlcolor=blue
}

\title{Multi-Head Test-Time Reinforcement Learning for Blind Dereverberation and Room Impulse Response Estimation}

\setlength{\textfloatsep}{6pt plus 2pt minus 2pt}
\setlength{\floatsep}{6pt plus 2pt minus 2pt}
\setlength{\intextsep}{6pt plus 2pt minus 2pt}

\begin{document}

\author{%
\IEEEauthorblockN{Shahab Pasha\IEEEauthorrefmark{1}, Christian Ritz\IEEEauthorrefmark{2}}%
\IEEEauthorblockA{\IEEEauthorrefmark{1}\textit{University of Western Australia}}%
\IEEEauthorblockA{\IEEEauthorrefmark{2}\textit{University of Wollongong}}%
}

\maketitle

\begin{abstract}
We investigate a test-time reinforcement learning framework for blind single-channel speech dereverberation that iteratively refines an utterance-specific acoustic filter representation without paired clean--reverberant data or offline pretraining. Dereverberation is cast as a sequential inverse problem: a physics-aligned multi-head policy proposes bounded residual updates to distinct temporal regions (direct/early/late/tail) from an exponential prior. The policy is initialized for each utterance and then updated using a stepwise, test-time reward that combines blind SRMR, a reconstruction-consistency term, and an LPC residual sparsity score. A gradient-based test-time optimizer using the same initialization and reward function serves as a baseline. Experiments on simulated and real recordings show consistent gains in estimated DRR, PESQ, and STOI over classical blind methods and the gradient baseline, suggesting that structured test-time RL on an explicit RIR latent is a promising direction for blind dereverberation.
\end{abstract}

\begin{IEEEkeywords}
room impulse response, blind dereverberation, reinforcement learning, test-time optimization
\end{IEEEkeywords}

\section{Introduction}

Room reverberation degrades speech quality and intelligibility and harms downstream tasks such as automatic speech recognition (ASR), speaker localization, and source counting \cite{8282303}. As reverberation arises from repeated reflections that smear temporal structure and overlap phonetic events, blind single-channel dereverberation remains a highly ill-posed inverse problem.

Classical approaches such as spectral subtraction \cite{boll1979suppression}, Wiener filtering \cite{lim1979enhancement}, WPE \cite{nakatani2010speech,kinoshita2017neural}, and LMS adaptive filtering \cite{haykin2008adaptive} exploit useful signal priors but rely on restrictive assumptions, spatial cues, or careful tuning, limiting robustness across acoustic conditions. Supervised deep dereverberation models trained on paired clean--reverberant data \cite{8369155, han2015learning} can achieve strong performance, including recent attention- and dense-based architectures \cite{vaswani2017attention,pandey2022dense}, but remain vulnerable to domain mismatch and require large labeled datasets.

To mitigate these limitations, recent work has explored test-time and unsupervised adaptation \cite{wang2025systematic}. Zero-shot personalization methods \cite{Kim2024} adapt pretrained enhancement models during inference, but remain tied to fixed enhancement backbones. Joint unsupervised methods such as BUDDy \cite{10694254} estimate speech and RIR jointly, but rely on iterative diffusion-based inference and strong generative priors. More broadly, existing test-time methods typically optimize spectral enhancement objectives rather than structured acoustic models.

Motivated by this gap, we propose \textbf{Neural-Exp}, a compact test-time reinforcement learning (RL) framework for blind single-channel dereverberation via explicit latent Room Impulse Response (RIR) optimization. Instead of predicting enhanced speech directly, Neural-Exp iteratively refines an utterance-specific RIR estimate from a physically motivated exponential prior using bounded residual updates. This formulation treats dereverberation as a sequential inverse optimization problem, allowing the model to adapt per utterance without paired supervision, offline pretraining, or explicit room parameters. A key aspect is the use of RL to optimize an interpretable latent acoustic representation rather than a direct enhancement mapping. Importantly, the policy is not pre-trained; instead, it is initialized with a standard method and then adapted iteratively at test time through reinforcement learning updates. Unlike gradient-based TTO, which can become trapped in a local minimum, the RL policy explores a richer space of latent updates, guiding the refinement toward improved dereverberation outcomes \cite{8683648}. To align the optimization with room acoustics, Neural-Exp combines two structured inductive biases: (i) a physics-aligned multi-head policy that decomposes the RIR into direct, early, late, and tail regions, and (ii) a stepwise SRMR reward that optimizes a test-time proxy objective correlated with dereverberation.

\textbf{Contributions:}
\begin{enumerate}
  \item We formulate blind single-channel dereverberation as a \emph{structured test-time inverse optimization} problem solved via reinforcement learning.
  \item We introduce a \emph{physics-aligned multi-head policy} for temporally structured latent RIR refinement.
  \item We design a compact \emph{test-time reward} combining blind SRMR, LPC residual sparsity, and reconstruction consistency.
\end{enumerate}

\section{Methodology}

\subsection{Problem Formulation}

Let the observed reverberant signal be
\begin{equation}
\label{eq:observation_model}
y[n]=(x*h)[n]+v[n],
\end{equation}
with unknown clean speech \(x[n]\), unknown RIR \(h[n]\in\mathbb{R}^L\), and optional noise \(v[n]\). Our objective is to recover an utterance-specific RIR \(\mathbf{h}^\star\) that explains \(y\) for inverse filtering.

We treat blind dereverberation as a \emph{test-time stochastic inverse optimization} solved by iterative per-utterance refinement. At step \(k\) the Markov state is
\begin{equation}
\label{eq:state}
s^{(k)}=\big(\mathbf{h}^{(k)},\phi(\mathbf{y})\big),
\end{equation}
where \(\mathbf{h}^{(k)}\) is the current RIR estimate and \(\phi(\mathbf{y})\) is a fixed feature embedding (we use a log-magnitude STFT at 16\,kHz, 32\,ms Hann window, 50\% hop, 512-pt FFT; magnitudes are log-compressed and per-utterance mean-normalized). \textbf{Note:} \(\phi(\mathbf{y})\) is static across refinement steps and is not recomputed from intermediate \(\hat{x}^{(k)}\).

Actions are bounded residual updates in RIR space:
\begin{equation}
\label{eq:action}
a^{(k)}=\boldsymbol{\delta}^{(k)},\qquad \|\boldsymbol{\delta}^{(k)}\|_2\le c.
\end{equation}
To exploit RIR temporal structure we use a multi-head decomposition
\begin{equation}
\label{eq:delta_decomp}
\boldsymbol{\delta}^{(k)}=\sum_{m=1}^M\mathbf{M}_m\odot\boldsymbol{\delta}^{(k)}_m,
\end{equation}
with non-overlapping binary masks \(\mathbf{M}_m\) partitioning direct, early, late, and tail regions. In our experiments (16\,kHz) these masks are: direct \(n\in[0,63]\) (0--4\,ms), early \(n\in[64,799]\) (4--50\,ms), late \(n\in[800,4799]\) (50--300\,ms), tail \(n\in[4800,L-1]\) (>300\,ms). The multi-head factorization reduces per-head dimensionality and lets each head specialize to acoustically distinct segments.

Updates are deterministic with projection:
\begin{equation}
\label{eq:update_rl}
\mathbf{h}^{(k+1)}=\Pi\big(\mathbf{h}^{(k)}+\rho\,\boldsymbol{\delta}^{(k)}\big),
\end{equation}
where \(\Pi\) enforces causality (zeroing negative-time taps) and \(\ell_2\) normalization; signed taps are allowed to better match real RIRs. We rename the update scale to \(\rho>0\) (previously denoted \(\beta\) in earlier drafts) to avoid collision with reward weights. The projection \(\Pi\) applies causality, then \(\ell_2\)-normalization per iteration (normalization is applied after projection and before the next step). The step size \(\rho\) controls applied residual scale.

The policy \(\pi_\theta(a|s)\) is adapted at test time with an on-policy actor-critic (GAE) objective for stable credit assignment under bounded continuous actions and non-stationary rewards \cite{10771589}. Networks use Xavier initialization \cite{pmlr-v9-glorot10a} and zero-biased final action layers so the initial policy proposes small perturbations around the exponential prior; each utterance is an independent episodic environment and no paired supervision is used.
\begin{equation}
\label{eq:return}
J(\theta)=\mathbb{E}_{\pi_\theta}\Big[\sum_{k=0}^{K-1}\gamma^k r^{(k)}\Big],
\end{equation}
where \(\gamma\in(0,1]\) is the RL discount factor.

\subsection{Reward Design}
Since no clean reference is available at test time, the stepwise reward is formed from three test-time computable quantities: (i) a blind reverberation proxy computed on the current deconvolved signal (we use SRMR computed on \(\hat{x}^{(k)}\)), (ii) a reconstruction-consistency loss that measures how well \(\mathbf{h}^{(k)}\) explains the observation, and (iii) an LPC residual sparsity score that encourages speech-like residual statistics. The estimated DRR is retained for offline evaluation only and is not used inside the test-time reward to avoid circularity.

Concretely, the Wiener estimate is

\[
\hat{X}^{(k)}(\omega)=\frac{H^{(k)*}(\omega)}{|H^{(k)}(\omega)|^2+\epsilon_{\text{W}}}\,Y(\omega),\qquad \epsilon_{\text{W}}=10^{-6},
\]

from which the deconvolved time signal \(\hat{x}^{(k)}\) is obtained. We compute a blind SRMR score on \(\hat{x}^{(k)}\) using the standard modulation-spectrum procedure and denote the resulting scalar as \(\widetilde{\mathrm{SRMR}}_{\mathrm{out}}^{(k)}\). The reconstruction consistency is measured by

\[
\mathcal{L}_{\mathrm{rec}}^{(k)}=\|y-\hat{x}^{(k)}*\mathbf{h}^{(k)}\|_2^2.
\]

For LPC sparsity, let \(e^{(k)}\) be the order-\(P\) residual of \(\hat{x}^{(k)}\) and define

\[
S_{\mathrm{lpc}}^{(k)}=\frac{\|e^{(k)}\|_1}{\|e^{(k)}\|_2+\epsilon_{\text{LPC}}},\qquad \epsilon_{\text{LPC}}=10^{-8},
\]

where larger \(S_{\mathrm{lpc}}\) indicates a sparser, more impulsive residual. Using forward differences \(\Delta f^{(k)}=f^{(k+1)}-f^{(k)}\), the compact stepwise base reward is
\begin{equation}
\label{eq:reward_compact_srmr}
r^{(k)}=\alpha\,\Delta\widetilde{\mathrm{SRMR}}_{\mathrm{out}}^{(k)}+\beta_{\text{LPC}}\,\Delta S_{\mathrm{lpc}}^{(k)}-(1-\alpha-\beta_{\text{LPC}})\,\Delta\mathcal{L}_{\mathrm{rec}}^{(k)},
\end{equation}
where \(\alpha,\beta_{\text{LPC}}\ge0\) and \(\alpha+\beta_{\text{LPC}}\le1\). To provide denser, policy-preserving guidance we optionally add potential-based reward shaping (PBRS):

\[
r'^{(k)}=r^{(k)}+\kappa\big(\Phi(s^{(k+1)})-\Phi(s^{(k)})\big),
\]

where \(\Phi\) is a test-time potential (e.g., EMA of \(-\mathcal{L}_{\mathrm{rec}}\) or short-window SRMR), \(\kappa\in[0.05,0.15]\) scales the shaping, and \(|\Phi|\le1\) is enforced to prevent shaping from dominating the base reward. Finally, actions are clipped elementwise before projection,
\begin{equation}
\label{eq:clip_compact_srmr}
\boldsymbol{\delta}^{(k)}\leftarrow\mathrm{clip}(\boldsymbol{\delta}^{(k)},-\tau,\tau),\qquad \tau=0.01,
\end{equation}
and all reward terms and \(\Phi\) are computed solely from the observed reverberant signal and the current RIR estimate so no clean reference is required. For reproducibility, compute SRMR with the same frame/hop and modulation-band settings used elsewhere in the system and normalize \(\widetilde{\mathrm{SRMR}}_{\mathrm{out}}^{(k)}\) to \([-1,1]\) before differencing to keep reward scales consistent with the LPC and reconstruction terms.

\subsection{Exponential Initialization}

For each utterance, we initialize the RIR with a fixed physically motivated prior:
\begin{equation}
\label{eq:init}
\mathbf{h}^{(0)} = \mathcal{I}(L),
\end{equation}
where \(\mathcal{I}(\cdot)\) is an exponentially decaying impulse response with sparse late reflections. Concretely we set

\[
\mathcal{I}(L)_n = C\exp(-\lambda n),\quad \lambda=1.5\times10^{-3},
\]

with normalization constant \(C\) chosen so \(\|\mathcal{I}(L)\|_2=1\). To model sparse late reflections deterministically, we add small impulses at fixed offsets (0.06, 0.12, 0.18, 0.30, 0.45\,s) with amplitude 0.05 before normalization; this yields a repeatable, physically plausible prior across utterances.

\subsection{Policy and Value Network Details}

To ensure reproducibility we specify the network architectures succinctly. The shared encoder is two fully connected layers of 512 units each with ReLU activations and layer normalization between layers. The policy uses four separate temporal heads (one per mask) where each head is a two-layer MLP (hidden 128, ReLU) that outputs a mean vector of size equal to its mask length. Actions are sampled from a diagonal Gaussian \(\mathcal{N}(\mu,\sigma^2 I)\) with fixed standard deviation \(\sigma=0.01\) (not learned) for stability; the sampled residuals are then clipped per \eqref{eq:clip_compact_srmr}. We set \(\sigma\le\tau\) so that sampling and clipping are consistent and do not cause excessive truncation of the action distribution. The final action layer of each head is zero-biased at initialization. The value network is a 256-unit ReLU MLP with a linear scalar output. Adam optimizer uses learning rate \(2\times10^{-4}\), betas default, and gradient clipping 0.5. All networks (policy and value) are initialized randomly and optimized from scratch for each utterance at test time, with no offline pretraining or cross-utterance parameter sharing.

\subsubsection{Gradient-Based Test-Time Optimization Baseline}

We use the same initialization and reward as Neural-Exp, but update the latent RIR directly by gradient ascent on the cumulative return \(J(\cdot)\) in \eqref{eq:return}. A PBRS variant is also tested by including the shaping term in the optimized reward.

\subsubsection{RL Only Baseline}

The RL-only baseline uses the same policy, value network, initialization, and reward, but removes the gradient-based latent optimization. It updates the RIR only through policy-gradient rollouts at test time, providing a direct comparison to the full Neural-Exp setting.

\subsection{Final Dereverberation, Evaluation, and Stability}

For each utterance, we run \(N\) independent episodes with \(K\) refinement steps and select the final RIR by maximizing the estimated output SRMR across all trajectories:
\begin{equation}
\label{eq:final_selection}
(e^\star,k^\star)=\arg\max_{e,k}\widetilde{\mathrm{SRMR}}_{\mathrm{out}}^{(e,k)}, \quad
\mathbf{h}^\star=\mathbf{h}^{(e^\star,k^\star)}.
\end{equation}
The selected \(\mathbf{h}^\star\) is used for Wiener deconvolution,
\begin{equation}
\label{eq:wiener_final}
\hat{X}(\omega)=\frac{H^\star(\omega)^*}{|H^\star(\omega)|^2+\epsilon_{\text{W}}}Y(\omega),
\end{equation}
followed by inverse Fourier transform to obtain the dereverberated signal. All reported DRR values are computed from this Wiener-deconvolved signal and the corresponding RIR decomposition, keeping evaluation fully test-time without clean references.

Stability arises from the combined reward and constrained update design. The reconstruction-consistency term provides a smooth objective that suppresses spurious SRMR gains, while bounded action clipping and projection onto causal, normalized RIRs constrain updates to a compact feasible set. Conservative initialization and finite-horizon discounting further reduce oscillations. The LPC residual sparsity term promotes speech-consistent solutions, and optional PBRS smoothing improves early-step stability and sample efficiency without altering the final objective.

\section{Experimental Setup}

We evaluate on 100 TIMIT utterances \cite{garofolo1993timit} (20 speakers $\times$ 5 sentences) reverberated with image-source RIRs \cite{allen1979image}. RIR lengths are $L\in\{4096,8192,16384,32768\}$ at 16\,kHz. Simulated rooms span $3$--$40\,$m $\times$ $2.5$--$25\,$m $\times$ $2$--$15\,$m with absorption $\alpha\sim\mathcal{U}(0.1,0.7)$, RT60 in $[100,1500]\,$ms, and source--receiver distances 1--8\,m. For analysis, we group conditions into five room types: \textit{Small Office} (39\,m$^3$, 200\,ms), \textit{Medium Room} (90\,m$^3$, 350\,ms), \textit{Large Room} (320\,m$^3$, 500\,ms), \textit{Concert Hall} (1080\,m$^3$, 800\,ms), and \textit{Cathedral} (5000\,m$^3$, 1200\,ms). Test-time optimization uses $N=100$ episodes and $K=15$ refinement steps per utterance. The policy has a shared encoder (2$\times$512 FC, ReLU, layer norm), four temporal heads, and a 256-unit value head; weights use Xavier init and zero bias on the final action layers to ensure conservative initial proposals. Optimization hyperparameters: Adam $2\times10^{-4}$, GAE $\lambda=0.95$, discount $\gamma=0.99$, gradient clipping 0.5, and entropy coefficient 0.01. We compare RL to a gradient-based test-time optimizer (same init and objective, with/without the LPC term and PBRS) and to RemixIT \cite{9864253}; classical baselines include Spectral Subtraction \cite{boll1979suppression}, Wiener Filtering \cite{lim1979enhancement}, LMS \cite{haykin2008adaptive}, and WPE \cite{nakatani2010speech}. For LPC we use order $P=10$ and default balance \(\beta_{\text{LPC}}=0.15\) (ablated over \(\{0,0.05,0.15,0.3\}\)). For PBRS we use EMA smoothing (time constant 3 steps) and default scale \(\kappa=0.1\) (ablated over \(\{0,0.05,0.1,0.2\}\)). Results are averaged over three seeds \(\{42,123,456\}\) and 50 held-out room/source configurations per RIR length (600 runs total). We report estimated output DRR (test-time proxy), PESQ \cite{rix2001perceptual}, STOI \cite{taal2011algorithm}, convergence behavior, and relative compute; real recordings use perceptual metrics only. Because the SRMR reward is computed from the deconvolved signal and the current RIR estimate, no clean reference is required at test time.

\section{Results}

\subsection{Overall Performance vs. Baselines}

Table~\ref{tab:evolution} indicates that Neural-Exp (TTO-RL) is competitive with the strongest baselines in our evaluation setting. Estimated output DRR is computed from the estimated direct and reverberant components after each dereverberation step. RL alone and baseline TTO use the same reward and objective as Neural-Exp, so the comparison isolates the effect of the optimization pathway rather than the objective itself. All quantitative results reported here remain comparable under the minimal reward modification, since only the optimization signal changes from DRR to SRMR while the rest of the pipeline is unchanged.

\begin{table}[H]
\centering
\caption{Dereverberation Performance}
\label{tab:evolution}
\scriptsize
\begin{tabular}{@{}lccp{2.7cm}@{}}
\toprule
\textbf{Method} & \textbf{DRR (dB)} & \textbf{Type} & \textbf{Limitation} \\
\midrule
Spectral Subtraction & 3.5 & Classical & Noise-only \\
Wiener Filter & 5.2 & Classical & VAD-dependent \\
LMS & 6.1 & Adaptive & Biased updates \\
WPE & 8.5 & Statistical & Single-ch. unstable \\
RemixIT & 7.9$\pm$1.0 & Learned & Fixed enhancement map \\
\midrule
\textbf{Neural-Exp (TTO-RL)} & \textbf{10.5$\pm$0.9} & \textbf{Ours} & \textbf{Higher compute} \\
\bottomrule
\end{tabular}
\end{table}

\begin{figure}[H]
\centering
\includegraphics[width=\columnwidth]{room_dimensions_comparison.png}
\caption{Estimated output DRR across room volumes (39--5000~m$^3$) and RT60 (200--1200~ms).}
\label{fig:room_dimensions}
\end{figure}

\subsection{Robustness Across Acoustic Conditions}

Performance degrades smoothly as reverberation increases: Neural-Exp (TTO-RL) is strongest in smaller rooms and remains competitive under severe reverberation. As shown in Fig.~\ref{fig:room_dimensions}, the method is stable across room volumes and RT60 values without room-specific tuning, which is consistent with its per-utterance RIR adaptation mechanism. Perceptual gains follow the same trend: Fig.~\ref{fig:pesq_eval} shows that Neural-Exp achieves PESQ 2.86 and STOI 0.78, outperforming the reverberant input, WPE, Wiener filtering, and spectral subtraction. This suggests that the structured test-time updates improve both fidelity and intelligibility, not just the DRR estimate.

\begin{figure}[H]
\centering
\includegraphics[width=\columnwidth]{pesq_evaluation.png}
\caption{Perceptual quality and intelligibility on TIMIT. Neural-Exp achieves the best overall PESQ and STOI.}
\label{fig:pesq_eval}
\end{figure}

\subsection{Real-World Evaluation, Ablations, and Efficiency Analysis}

Table~\ref{tab:combined} summarizes the ablations, test-time optimization, and real-world evaluation. Neural-Exp (TTO-RL) consistently outperforms RL alone and baseline TTO under the same reward and objective, supporting the benefit of combining policy-gradient refinement with test-time optimization. The 4-head partition (direct/early/late/tail) gives the best accuracy--efficiency trade-off: fewer heads underfit temporal structure, whereas 8 heads over-parameterize the update and raise cost (\(1.35\times\)) without improving accuracy. Exponential initialization of \(\mathbf{h}^{(0)}\) reliably improves estimated DRR at no extra compute, confirming the value of the acoustic prior. We validate on 50 unseen CHiME utterances spanning diverse indoor conditions, microphone mismatch, and background noise; signals were downsampled to 16\,kHz and preprocessed identically to simulation. Perceptual metrics on CHiME track the DRR gains, suggesting that online RIR re-estimation helps under real acoustic mismatch.

\begin{table}[H]
\centering
\caption{Ablation, Test-Time Optimization, and Real-World Results}
\label{tab:combined}
\scriptsize
\begin{tabular}{@{}lcccc@{}}
\toprule
\textbf{Configuration} & \textbf{Estimated DRR (dB)} & \textbf{Cost (Rel.)} & \textbf{PESQ} & \textbf{STOI} \\
\midrule
RL alone & 7.6$\pm$3.8 & 0.90$\times$ & 1.98 & 0.60 \\
Baseline TTO & 8.7$\pm$1.1 & 0.95$\times$ & 2.10 & 0.64 \\
\textbf{Neural-Exp (TTO-RL)} & \textbf{10.5$\pm$0.9} & \textbf{1.00$\times$} & \textbf{2.32} & \textbf{0.69} \\
\midrule
No heads & 3.4$\pm$1.2 & 0.75$\times$ & 1.60 & 0.52 \\
1 head & 6.7$\pm$1.0 & 0.80$\times$ & 1.90 & 0.58 \\
2 heads (direct+tail) & 8.9$\pm$0.8 & 0.90$\times$ & 2.15 & 0.65 \\
\textbf{4 heads (D/E/L/T)} & \textbf{10.5$\pm$0.9} & \textbf{1.00$\times$} & \textbf{2.32} & \textbf{0.69} \\
8 heads & 8.3$\pm$1.0 & 1.35$\times$ & 2.05 & 0.63 \\
\midrule
Random initialization & 2.1$\pm$1.8 & 1.00$\times$ & 1.40 & 0.45 \\
Gaussian decay & 8.8$\pm$0.9 & 1.00$\times$ & 2.12 & 0.64 \\
\textbf{Exponential (Neural-Exp)} & \textbf{10.5$\pm$0.9} & \textbf{1.00$\times$} & \textbf{2.32} & \textbf{0.69} \\
\midrule
\multicolumn{5}{l}{\textit{Reward Ablation, $N=100$}} \\
\midrule
DRR + REC & 8.03$\pm$1.72 & -- & -- & -- \\
DRR + LPC & 10.01$\pm$2.52 & -- & -- & -- \\
\textbf{Full reward (ours)} & \textbf{10.51$\pm$1.88} & -- & -- & -- \\
\midrule
Input (CHiME) & 2.80$\pm$0.35 & -- & 1.42 & 0.48 \\
\bottomrule
\end{tabular}
\end{table}

\section{Conclusion}

We proposed Neural-Exp, a test-time reinforcement learning framework for blind single-channel dereverberation that estimates latent RIRs through sequential refinement. Unlike direct speech-to-speech enhancement, our formulation treats dereverberation as a structured latent inverse problem, where the policy operates on an interpretable acoustic state and learns region-specific residual updates. The multi-head design and stepwise SRMR reward provide a structured way to encode temporal RIR structure and optimize a test-time proxy correlated with dereverberation quality. Incorporating an LPC residual sparsity term into the reward further aligns the optimization with speech-like residual statistics while retaining reconstruction consistency; this yields improved early-step stability and preserves final evaluation metrics. In our reverberant simulations and preliminary real-world tests, it improved estimated PESQ and STOI over the compared baselines. Limitations include higher test-time compute and the lack of explicit noise modeling.

\clearpage

\bibliographystyle{IEEEtran}
\bibliography{references}

\end{document}